\scanpage{27}
\mspage{25}
% archive sheet number: 26
% hand: Grothendieck
% sheet_kind: text
% element: text | top, line 1 | Si
Si
% element: display | top, line 2 | N=Ker(u:pi to G(R))
\[
N=\operatorname{Ker}(u:\pi\longrightarrow G(\unsure{R|\mathbb Q_\ell})\subset\operatorname{Aut}(E)),\quad(\inserted{\text{donc }}N\subset U),
\tag*{(5.6.1)}
\]
% element: text | upper, line 3 | on voit donc que
on voit donc que
% element: display | upper, line 4 | U/U cap N into G^0(R)
\[
\deleted{\unsure{E}}U/U\cap N\hookrightarrow G^\circ(R)
\]
% element: text | upper-middle, lines 5–8 | s'identifie à un Z_l-module
s'identifie à un $\mathbb Z_\ell$-module libre de rang $d$ qui engendre $G^\circ(R)$ comme vectoriel sur $\mathbb Q_\ell$. D'ailleurs, on trouve ainsi un hom d'extensions
% element: diagram | middle, lines 9–10 | two exact sequences
\diagram{\begin{tikzcd}[ampersand replacement=\&,column sep=small]
1\arrow[r]\&U/N\arrow[r]\arrow[d,hook]\&\pi/N\arrow[r]\arrow[d,hook]\&\pi/U\arrow[r]\arrow[d,"\sim"]\&1\\
1\arrow[r]\&G^\circ(R)\arrow[r]\&G(R)\arrow[r]\&G(R)/G^\circ(R)\arrow[r,dash]\&1
\end{tikzcd}}
% element: text | middle-lower, lines 11–17 | Ainsi
Ainsi, l'\deleted{\unsure{extension}} groupe algébrique $G(R)$ est \deleted{\unsure{canoniq.}} canoniquement déterminé (à isom unique près) par la connaissance du sous-groupe \inserted{(ouvert)} $U$ de $\pi$, et du sous-groupe \inserted{fermé} $N$ de $\pi$ tel que $U/N\simeq\mathbb Z_\ell^d$. $G(R)$ est en effet l'image de l'extension canonique de $\pi/U$ par $U/N$ \deleted{\unsure{$G^\circ(R)$}} par l'hom canonique
% element: display | lower, line 18 | U/N to tensor
\[
U/N\unsure{-|\longrightarrow}U/N\otimes_{\mathbb Z_\ell}\mathbb Q_\ell
\]
% element: text | bottom, lines 19–22 | NB: la structure
[ NB: la structure de $\mathbb Z_\ell$-module sur le groupe profini $U/N$ est canoniquement déterminée), et la structure de groupe algébrique se détermine aussi trivialement .....
% element: text | bottom edge, lines 23–25 | Si M_l
Si $\unsure{{}'M_\ell|M_\ell|M'_\ell}=\underbrace{U/[U,U]^{(\ell)}}_{\text{\unsure{compl|compl.} $\ell$-primaire}}$ est un $\mathbb Z_\ell$-module de t.f. alors on construit une extension \uline{universelle} de $\Gamma=\pi/U$ par le groupe vectoriel correspondant à $M_\ell\otimes_{\mathbb Z_\ell}\mathbb Q_\ell$.
